
\def \Subject {طبقه بندی با کمک شبکه های مدرن پیچشی }
\def \Session { دو}
% \setcounter{chapter}{\Session-1}
\renewcommand{\chaptername}{سوال}
\chapter{\Subject}

\section{ ساخت دیتالودر به صورت دستی}

 اول برای مجموعه داده، 
ابعاد آن به ۲۸ در ۲۸ تغییر داده شد تا لایه‌های پیچشی دو بعدی 
\LTRfootnote{2d convolutional}
بتوانند کار کنند. برای افزونگی داده، سه تغییر چرخش، انتقال و 
نویز نمک‌فلفلی
\LTRfootnote{Salt and pepper noise}
صرفا روی مجموعه‌داده آموزش با ضریب احتمال $0.2$ اضافه شدند. از آنجایی که حجم مجموعه‌داده بزرگ است، برای این که زمان آموزش خیلی زیاد نشود، بخش افزونگی داده تصویر جدیدی درست نمی‌کند و صرفا تصویر قبلی را تغییر می‌دهد. همچنین بر روی همه مجموعه‌ها (آموزش، صحت‌سنجی و آزمون) چک می‌شود که اگر مقدار شدت رنگ بصورت ۰ تا ۲۵۵ داده شده، به صفر تا یک مقیاس 
\LTRfootnote{Scale}
شود. در بخش 
 دیتالودر
\LTRfootnote{DataLoader}
مقدار ۳۲ برای سایز دسته
\LTRfootnote{Batch size}
درنظر گرفته شد و تعیین شد تا داده‌‌ها بصورت تصادفی جابجا شوند.
\LTRfootnote{Shuffle}


\section{
جایگذاری بلوک 
\lr{A}
}

در این بخش از 
بلوک 
\lr{A}
استفاده شد. نکته خاص بلوک 
\lr{A}
این است که با اضافه کردن یک اتصال باقی‌مانده 
\LTRfootnote{Residual connection}
باعث می‌شود تا بلوک ویژگی‌های ورودی را از دست ندهد و هم ویژگی‌های ورودی و هم ویژگی‌های استخراجی در خروجی آن نقش دارد.

تفاوت دو نوع معماری داده شده در این است که یکی از آن‌ها در ارتباط مستقیم ورودی، یک لایه پیچش یک در یک دارد. این باعث می‌شود تا این بلوک مجبور نباشد تا تعداد کانال‌های ورودی و خروجی آن برابر باشد، زیرا اگر این لایه نبود، هنگام جمع ورودی و خروجی بلوک، ابعاد متفاوت می‌شد ولی با وجود این بلوک، ارتباط مستقیم از یک لایه کانولوشن عبور می‌کند که کانال‌های ورودی آن برابر با کانال‌های ورودی بلوک و کانال‌های خروجی آن برابر با کانال‌های خروجی بلوک است. 

این لایه یک در یک علاوه بر تطبیق ابعاد، این امکان را نیز به مدل می‌دهد تا عمق بیشتری (کانال‌های بیشتری) ‫‫‫از تصویر را پردازش کند و سپس این ویژگی‌ها را باهم ترکیب کند. 

        \begin{figure}[H]
	\centering
	\includegraphics[width=\textwidth]{images/q2a.jpg}
	\caption{
		شاخص‌های مختلف مدل \lr{A} هنگام آموزش
	}
	\label{fig:q2a}
\end{figure}


        \begin{table}[H]
	\centering
	\caption{ شاخص‌های مختلف مدل \lr{A} هنگام آموزش }
	\begin{latin}


\begin{tabular}{lrrrr}
	\toprule
	Epoch & train\_loss & train\_acc & val\_loss & val\_acc \\
	\midrule
	0 & 0.155060 & 95.310000 & 0.058931 & 98.360000 \\
	1 & 0.055143 & 98.426000 & 0.025974 & 99.240000 \\
	2 & 0.046224 & 98.662000 & 0.023006 & 99.380000 \\
	3 & 0.039257 & 98.890000 & 0.020706 & 99.430000 \\
	4 & 0.035947 & 98.968000 & 0.026350 & 99.200000 \\
	\bottomrule
\end{tabular}
	\end{latin}
\end{table}

\section{جایگذاری بلوک \lr{B}}
ایده بلوک \lr{B} که از مدل اینسپشن 
\LTRfootnote{Inception}
الهام گرفته شده‌است، این است که در شبکه‌های پیچش، یکی از چالش‌ها تعیین اندازه هسته پیچش
\LTRfootnote{Convolution Kernel}
است که می‌تواند باعث شود ویژگی‌های خوبی استخراج نشود و یا آموزش مدل طولانی شود. این بلوک، هسته‌های مختلفی را امتحان کرده و درنهایت همه آن‌ها باهم خروجی را تشکیل می‌دهند.


\begin{figure}[H]
	\centering
	\includegraphics[width=\textwidth]{images/q2b.jpg}
	\caption{
		شاخص‌های مختلف مدل \lr{B} هنگام آموزش
	}
	\label{fig:q2b}
\end{figure}


\begin{table}[H]
	\centering
	\caption{ شاخص‌های مختلف مدل \lr{B} هنگام آموزش }
	\begin{latin}
\begin{tabular}{lrrrr}
	\toprule
	Epoch & train\_loss & train\_acc & val\_loss & val\_acc \\
	\midrule
	0 & 0.426373 & 85.952000 & 0.157498 & 95.360000 \\
	1 & 0.130969 & 96.054000 & 0.078504 & 97.440000 \\
	2 & 0.103298 & 96.940000 & 0.061154 & 98.210000 \\
	3 & 0.092428 & 97.144000 & 0.046544 & 98.730000 \\
	4 & 0.086451 & 97.344000 & 0.045756 & 98.740000 \\
	\bottomrule
\end{tabular}
	\end{latin}
\end{table}

\section{جایگذاری بلوک \lr{C}}
بلوک 
\lr{C}
عملا ترکیب دو بلوک قبل است. به این صورت که اتصال باقی‌مانده که در بلوک 
\lr{A}
بود را دارد و همچنین شبکه های پیچشی متعدد که در بلوک 
\lr{B}
موجود بود را نیز دارد.

پارامتر 
\lr{g}
عملا تعداد گروه‌های پیچشی را نشان می‌دهد (که برای مثال در بلوک 
\lr{B}
دو لایه پیچشی سه در سه وجود دارد. 
)

هرچقدر این گروه‌ها بیشتر باشند، اطلاعات متنوع‌تری (و نه لزوما عمیق‌تر) از داده‌ها یافت می‌شود ولی مدل سنگین‌تر می‌شود.
\begin{figure}[H]
	\centering
	\includegraphics[width=\textwidth]{images/q2c.jpg}
	\caption{
		شاخص‌های مختلف مدل \lr{B} هنگام آموزش
	}
	\label{fig:q2c}
\end{figure}


\begin{table}[H]
	\centering
	\caption{ شاخص‌های مختلف مدل \lr{C} هنگام آموزش }
	\begin{latin}
		
\begin{tabular}{lrrrr}
	\toprule
	Epoch & train\_loss & train\_acc & val\_loss & val\_acc \\
	\midrule
	0 & 0.486524 & 83.976000 & 0.182243 & 94.850000 \\
	1 & 0.162272 & 95.110000 & 0.062752 & 98.140000 \\
	2 & 0.123943 & 96.212000 & 0.063555 & 98.000000 \\
	3 & 0.103672 & 96.802000 & 0.075380 & 97.740000 \\
	4 & 0.095231 & 97.052000 & 0.098605 & 97.420000 \\
	\bottomrule
\end{tabular}
	\end{latin}
\end{table}

\section{نتایج}

\begin{figure}[H]
	\centering
	\includegraphics[width=\textwidth]{images/q2all.jpg}
	\caption{
		نمودار خطا و دقت هر سه مدل در کنار هم
	}
	\label{fig:q2all}
\end{figure}


\begin{table}[H]
	\centering
	\caption{ نتیجه نهایی مدل‌ها}

		
		\begin{tabular}{r c c c}
			\toprule
		مدل &	آموزش & صحت‌سنجی & آزمون \\
			\midrule

 تمرین اول &	$99.124$ & $97.5$ & $97.43$  \\
 بلوک \lr{A} & $99.27$ & $99.2$ & $98.068$ \\
بلوک \lr{B}  & $98.72$ & $98.74$ & $96.512$ \\	
بلوک \lr{C}  & $97.07$ & $97.42$ & $97.042$ \\
			\bottomrule
		\end{tabular}
\label{q2:comp}
\end{table}
همانطور که در نمودار 
\ref{fig:q2all}
می‌توان دید، بعد از پنج دور آموزش، همچنان خطای آموزش نزولی هست ولی خطای صحت‌سنجی دیگر نزولی نیست و مدل اشباع شده است. برای جلوگیری از بیش‌برازش از افزودن نویز به تصویر در بخش افزونگی داده و اضافه کردن لایه دراپ‌اوت 
\LTRfootnote{Dropout}
استفاده شد و خطای صحت‌سنجی به آموزش نزدیک شد ولی همچنان جای کار دارد.

با دیدن جدول 
\ref{q2:comp}
می‌توان دید که به غیر از مدل 
\lr{A}
بقیه مدل‌ها از مدل تمرین اول عملکرد ضعیف‌تری دارند. این موضوع به این دلیل است که مجموعه‌داده 
\lr{MNIST}
مجموعه‌داده به نسبت ساده‌ای است و یک مدل عصبی ساده ولی با بهینه‌سازی‌های جانبی (برای مثال استفاده از منظم‌سازی و تنظیم دقیق تعداد نورون و لایه) می‌تواند به خوبی آن را فرا بگیرد. ولی مدل‌های پیچیده‌تری مثل مدل‌های این سوال چون قدرت بیشتری دارند، احتمال دچار شدن به بیش‌برازش نیز در آن‌ها بیشتر است. با افزایش تعداد دورهای یادگیری به ۵۰ (تعداد دوره آموزشی در تمرین اول) می‌توان با اطمینان خوبی حدس زد که خطای آموزش در این مدل‌ها بسیار کمتر باشد ولی خطای صحت‌سنجی بهبودی نیابد و یا حتی بدتر شود. 

تنها مدلی که از مدل تمرین قبل بهتر عمل کرد مدل با بلوک 
\lr{A}
است. زیرا با استفاده از اتصال باقی‌مانده، این مدل هم ویژگی‌های ساده (ویژگی‌های ورودی) و هم ویژگی‌های مکانی پیچیده‌تری که با لایه پیچشی یاد گرفته است را حفظ می‌کند و درنتیجه در این تسک ساده، داده‌ای را از دست نمی‌دهد و اطلاعات بیشتری را برای لایه‌های تماما متصل
\LTRfootnote{Fully connected}
فراهم می‌کند. البته حتی این موضوع هم بعد از افزودن نویز و لایه دراپ‌اوت ممکن شد.

به عنوان مثال. بلوک 
\lr{C}
که به عنوان ترکیب دو بلوک دیگر است از هردوی آن‌ها ضعیف‌تر عمل کرده است. روند مشهود در آموزش به این صورت است که هرچقدر مدل‌ها پیچیده‌تر شدند دقت کمتر شده است. پس عملا در اینجا کم‌برازش 
\LTRfootnote{Underfitting}
رخ نداده است و اصلا از اول نیز نیازی به استفاده از مدل‌های پیچیده نبوده.

\section{نمایش خروجی فیلترهای کانولوشنی}

\begin{figure}[H]
	\centering
	\includegraphics[width=0.5\textwidth]{images/q2first_layer_orig.jpg}


	\centering
	\includegraphics[width=\textwidth]{images/q2first_layer_chans.jpg}
	\caption{
یک عکس ورودی و خروجی فیلترهای لایه‌ اول مدل برای آن
	}
	\label{fig:q2_first_layer}
\end{figure}
با توجه به عکس 
\ref{fig:q2_first_layer}
می‌توان دید که فیلترها در اینجا کارهایی مثل لبه‌یابی عمودی (مثلا کانال ۶)، لبه‌یابی افقی (کانال ۹) و یا وارونه کردن رنگ‌های تصویر (کانال ۲) را انجام داده‌اند.  فیلترهای ساده‌ای که در بینایی ماشین کلاسیک بصورت دستی تهیه و به مدل‌ها داده می‌شد.

\section{انتقال دانش}
با انجام این کار آموزش بسیار سریعتر شد. البته فریز کردن لایه در این‌جا خیلی منطقی نیست زیرا از مجموعه داده جدید نیز حجم خوبی موجود است. نتایج این آموزش در جدول \ref{q2:transfer_learn} آمده است. 

در بررسی‌هایی که انجام شد، فریز نکردن لایه‌ها باعث شد تا مطابق انتظار فرآیند آموزش طول بکشد ولی بعد از ۲۰ دوره، دقتی حدود ۷۰ درصد را نتیجه داد که از نتیجه‌های بدست آمده بهتر است.

\begin{table}[H]
	\centering
	\caption{ دقت مدل روی مجموعه داده آزمون \lr{fashion mnist}}
	
	
	\begin{tabular}{r c}
		\toprule
		مدل &	دقت  \\
		\midrule
		

		بلوک \lr{A}  & 	$62.86$ \\
		بلوک \lr{B}  & 	$50.0$ \\
		بلوک \lr{C}  & 	$48.43$ \\
		\bottomrule
	\end{tabular}
	\label{q2:transfer_learn}
\end{table}

\begin{figure}[H]
	
	\centering
	\includegraphics[width=\textwidth]{images/q2all_fashion.jpg}
	\caption{
خروجی و دقت روی مجموعه‌داده‌های آموزش و صحت‌سنجی برای 
\lr{fashion mnist}
	}
	\label{fig:q2_transfer}
\end{figure}

\begin{figure}[H]
	
	\centering
	\includegraphics[width=\textwidth]{images/fashion_mnist_confusion_matrix.png}
	\caption{
ماتریس درهم‌ریختگی برای مدل با بلوک \lr{A} (بهترین مدل) روی مجموعه‌داده
		\lr{fashion mnist}
	}
	\label{fig:q2_confusion}
\end{figure}
همانطور که در تصویر 
\ref{fig:q2_confusion}
می‌توان دید، مدل خیلی خوب عمل نکرده‌است و بخش خوبی از تصاویر را در کلاس‌های صفر و چهار طبقه‌بندی کرده است (البته این موضوع ثابت نیست و در یک اجرای دیگر کلاس‌های شش و نه بیشترین عضو را داشتند.). دلیلی که برای این اشتباه می‌تواند باشد این است که ویژگی‌هایی که برای تشخیص اعداد بکار می‌رود عموما از جنس انحنای اعداد است ولی ویژگی‌هایی که برای طبقه‌بندی در مجموعه‌داده 
\lr{fashion mnist}
استفاده می‌شوند از جنس بافت لباس هستند. و برخلاف اعداد، لباس‌ها ممکن است انحنای یکسانی داشته باشند (برای مثال گروه صفر، دو، چهار و شش) و تنها تفاوت آن‌ها در جنس لباس باشد.

\begin{figure}[H]
	
	\centering
	\includegraphics[width=\textwidth]{images/fashion_mnist_samples.png}
	\caption{
		نمونه اعضای کلاس‌های مجموعه‌داده 
		\lr{Fashion mnist}
	}
	\label{q2:sample}
\end{figure}

اگر لایه استخراج ویژگی منجمد نشود،  ماتریس درهم‌رریختگی  
\ref{q2:nonfreeze_conf}
حاصل می‌شود که در آن دیگر مدل بایاس
\LTRfootnote{Bias}
روی گروه خاصی ندارد و خطا به نسبت پراکندگی دارد. گروه‌هایی که شباهت ظاهری دارند (برای مثال گروه ۲ و گروه ۴) باهم اشتباه گرفته می‌شوند.

\begin{figure}[H]
	
	\centering
	\includegraphics[width=\textwidth]{images/q2_nonfreeze.png}
	\caption{
ماتریس درهم‌ریختگی بدون منجمد کردن لایه استخراج ویژگی
	}
	\label{q2:nonfreeze_conf}
\end{figure}


